\chapter{Pipeline control}
\label{ch:pipectl}

\section{Pipeline registers}

Both pipeline registers --- IF/DX and DX/WB --- are built the same way: a
\sig{valid} bit plus a payload.

\sig{valid} is what makes a bubble safe. Nothing downstream looks at the payload
unless \sig{valid} is set, so a bubble is simply \sig{valid}~=~0 and a flush
produces a guaranteed no-operation whatever the payload happens to hold. The
alternative considered was injecting an \sig{addi x0,x0,0} encoding, which would
have required the flush path to drive 32 correct bits and every downstream
decoder to interpret them --- more hardware and more ways to be wrong. The
payload is loaded only when a real instruction moves, so bubbles and stalls
leave those flops quiet.

Out of reset the payload holds a defined idle state --- the canonical NOP
\texttt{0x0000\_0013} in the instruction word and \sig{RESET\_PC} in the
program-address fields --- so the decoder outputs that reach the AXI
\sig{VALID} signals are defined from the first clock.

\section{The hazard unit}

All stall, flush and forwarding decisions are made in one module. No stage
manages its own progress. The entire control is five equations
(\file{hdl/hazard\_unit.v}):

\begin{center}
\begin{minipage}{0.72\linewidth}
\ttfamily\small
s2\_advance\_o   = \textasciitilde lsu\_busy\_i \& \textasciitilde wfi\_wait\_i;\\
redirect\_o      = (trap\_take\_i | mret\_exec\_i | mispredict\_i) \& s2\_advance\_o;\\
if\_dx\_we\_o     = s2\_advance\_o;\\
if\_dx\_bubble\_o = redirect\_o | \textasciitilde fetch\_valid\_i;\\
dx\_squash\_o     = trap\_take\_i;
\end{minipage}
\end{center}

\subsection{Stall sources}

There are exactly two, and they are the two terms of \sig{s2\_advance}: a data
transaction in flight (\sig{lsu\_busy}), and a sleeping \sig{WFI}
(\sig{wfi\_wait}).

The feature specification lists three stall sources, including a load-use hazard
between S2 and S3 and an instruction fetch that is not ready. Neither is a stall
in this implementation, and the difference is deliberate:

\begin{itemize}[nosep]
  \item \textbf{Fetch not ready} does not stall S2. It produces a bubble
        (\sig{if\_dx\_bubble = redirect | \textasciitilde fetch\_valid}): the stage
        advances, simply without an instruction. Stalling would have been
        indistinguishable in throughput and would have added a term to
        \sig{s2\_advance} for no benefit.
  \item \textbf{Load-use} is not a hazard here at all. In a three-stage pipeline a
        load's data is already in S3 by the time its consumer enters S2, so a
        single S3$\to$S2 bypass covers it. The classic load-use stall never
        arises and costs zero cycles.
\end{itemize}

\subsection{Flush sources}

Three, exactly as specified: a misprediction resolved in S2, trap entry, and
\sig{MRET}. All three produce the same action --- assert \sig{redirect}, load a
bubble into IF/DX --- so there is one mechanism rather than three.

\subsection{Simultaneous stall and flush}

Flush wins, because the flush invalidates the very instruction the stall was
waiting for. There is a deliberate second half to the rule, and it comes from the
same gate: every flush term is ANDed with \sig{s2\_advance}, so a trap
\emph{cannot} leave S2 while a data transaction is still outstanding. The bus is
therefore never left with an orphaned transaction whose response nobody will
accept.

Both properties fall out of one signal. That is why the unit is five lines and
not a state machine.

\subsection{Forwarding}

The hazard unit also computes the forwarding hits: S3's destination register
compared against each of S2's source registers, with \sig{x0} excluded. A hit
substitutes the writeback value for the register file output.

\section{Decision priority in S2}

Every architectural decision is taken in S2, in a fixed sequential order:

\begin{center}
interrupt \; $>$ \; synchronous exception \; $>$ \; \sig{MRET} \; $>$ \; misprediction
\end{center}

The interrupt is tested first, and specifically before the instruction issues any
data transaction, so the preempted instruction can be re-run identically after
the handler returns. An exception outranks a misprediction because a faulting
instruction has no correct path left to take. The \sig{redirect\_pc} multiplexer
in \file{hdl/cpu\_top.v} is written in exactly this order.

\section{Timing of the common cases}

\begin{table}[H]
\centering
\caption{Cycle behaviour of the four cases that determine throughput}
\label{tab:cputiming}
\begin{tabular}{@{}lL{0.66\linewidth}@{}}
\toprule
\textbf{Case} & \textbf{Behaviour} \\
\midrule
ALU operation & One instruction enters S2 every cycle. Measured CPI is 1.00 over 33
  straight-line instructions on a latency-1 instruction memory. \\
Load / store  & S2 stalls for the whole AXI round trip. S1 may complete a fetch
  already started and park it in the holding register; nothing enters or leaves S2. \\
Misprediction & Costs exactly one cycle: S2 resolves the branch and asserts
  \sig{redirect}, the wrong-path fetch is dropped, S2 executes a bubble, and the
  correct-path instruction is in S2 two cycles after the branch. \\
Trap entry    & No extra cycle beyond the redirect. The offending instruction is
  squashed on its way into DX/WB in the same cycle the CSR state is committed. \\
\bottomrule
\end{tabular}
\end{table}

\section{Interrupt acceptance}

An interrupt is taken only when all of the following hold in the same cycle
(\file{hdl/cpu\_top.v}):

\begin{center}
\ttfamily\small
irq\_take = ifdx\_valid\_q \&\& !lsu\_active \&\& mie\_global \&\& irq\_deliverable
\end{center}

where \sig{irq\_deliverable} combines \sig{cpu\_irq} with the corresponding
\reg{mie} bit. The \sig{!lsu\_active} term is what holds an interrupt through a
multi-cycle AXI stall: the request waits for the response beat rather than
abandoning a transaction mid-flight. \sig{ifdx\_valid\_q} is what confines it to
an instruction boundary.

The claim and end-of-interrupt pulses are registered outputs, each exactly one
cycle wide:

\begin{itemize}[nosep]
  \item \sig{cpu\_irq\_ack} $\leftarrow$ \sig{irq\_take \&\& s2\_advance}
  \item \sig{cpu\_irq\_eoi} $\leftarrow$ \sig{mret\_exec \&\& s2\_advance \&\& in\_irq}
\end{itemize}

\sig{in\_irq} is an internal flag set when an interrupt is taken and cleared on
\sig{MRET}. It exists so that an \sig{MRET} closing a \emph{synchronous} trap
does not emit an end-of-interrupt pulse --- the PIC must pop exactly one nesting
level per interrupt, and no levels for exceptions.

\begin{figure}[H]
\centering
\begin{tikztimingtable}[timing/slope=0,timing/coldist=2pt,xscale=1.05]
  clk           & 14{C} \\
  cpu\_irq      & 2L 12H \\
  irq\_take     & 4L 2H 8L \\
  cpu\_irq\_ack & 6L 2H 6L \\
  cpu\_in\_trap & 6L 8H \\
  cpu\_irq\_eoi & 12L 2H \\
  \extracode
  \begin{pgfonlayer}{background}
    \vertlines[help lines]{2,4,6,12}
  \end{pgfonlayer}
\end{tikztimingtable}
\caption{Interrupt handshake. \sig{cpu\_irq\_ack} is registered one cycle after the
CPU decides to take the interrupt; \sig{cpu\_irq\_eoi} follows the
interrupt-returning \sig{MRET}. Both are one cycle wide. \sig{cpu\_in\_trap}
spans entry to return and is not used by the PIC.}
\label{fig:irqhandshake}
\end{figure}

\section{\sig{WFI}}

\sig{WFI} is implemented as a third stall condition rather than as a
no-operation. S2 freezes on the instruction, S1 parks its fetch in the holding
register, and because no new AR is issued while that register is full the
instruction bus falls silent with no additional logic. Measurement in the
testbench shows at most the one fetch already in flight completing per sleep
window.

The wake condition is \sig{cpu\_irq \&\& mie[16+cpu\_irq\_vec]}. The global
\reg{mstatus.MIE} is deliberately not part of it, following the privileged
specification: a locally enabled interrupt resumes execution even when global
interrupts are disabled. \reg{MIE} then decides only what happens next --- with
\reg{MIE}~=~1 the wake becomes a normal interrupt, with \reg{mepc} = \sig{wfi+4}
so that \sig{MRET} resumes past the instruction rather than sleeping again; with
\reg{MIE}~=~0 the \sig{WFI} simply commits as a no-operation and execution falls
through.
